% !TEX program = xelatex
\documentclass[11pt,a4paper]{article}

\usepackage{fontspec}
\defaultfontfeatures{Ligatures=TeX,Scale=MatchLowercase}
\IfFontExistsTF{Times New Roman}{
  \setmainfont{Times New Roman}
}{
  \setmainfont{DejaVu Serif}
}
\IfFontExistsTF{Arial}{
  \setsansfont{Arial}
}{
  \setsansfont{DejaVu Sans}
}
\setmonofont{DejaVu Sans Mono}[Scale=0.82]

\usepackage[left=2cm,right=2cm,top=2cm,bottom=2cm,footskip=0.9cm]{geometry}
\usepackage{setspace}
\setstretch{1.15}
\usepackage{indentfirst}
\setlength{\parindent}{0cm}
\setlength{\parskip}{4.5pt}

\usepackage[table]{xcolor}
% Bảng màu chuyên nghiệp & hiện đại
\definecolor{PrimaryNavy}{HTML}{0F172A}
\definecolor{AccentBlue}{HTML}{1D4ED8}
\definecolor{SoftBlueBg}{HTML}{F8FAFC}
\definecolor{BorderBlue}{HTML}{CBD5E1}
\definecolor{TerminalBg}{HTML}{0F172A}
\definecolor{TerminalBorder}{HTML}{334155}
\definecolor{TerminalText}{HTML}{F8FAFC}
\definecolor{TerminalGreen}{HTML}{22C55E}
\definecolor{TerminalPrompt}{HTML}{38BDF8}
\definecolor{NoteBg}{HTML}{F0FDF4}
\definecolor{NoteBorder}{HTML}{86EFAC}
\definecolor{NoteText}{HTML}{166534}
\definecolor{WarnBg}{HTML}{FEF2F2}
\definecolor{WarnBorder}{HTML}{FECACA}
\definecolor{WarnText}{HTML}{991B1B}
\definecolor{TipBg}{HTML}{FFFBEB}
\definecolor{TipBorder}{HTML}{FDE68A}
\definecolor{TipText}{HTML}{92400E}
\definecolor{TableHeader}{HTML}{1E293B}
\definecolor{TableSubHeader}{HTML}{F8FAFC}
\definecolor{InlineCodeBg}{HTML}{F1F5F9}

\usepackage{graphicx}
\usepackage{tabularx}
\usepackage{booktabs}
\usepackage{enumitem}
\usepackage{listings}
\usepackage[most]{tcolorbox}
\usepackage{fancyhdr}
\usepackage{titlesec}
\usepackage{microtype}
\usepackage{hyperref}

\hypersetup{
  colorlinks=true,
  linkcolor=AccentBlue,
  urlcolor=AccentBlue,
  pdftitle={Hướng dẫn cài đặt và chạy dự án StayGo trên Windows},
  pdfauthor={Hệ thống Đặt phòng Khách sạn StayGo},
  pdfsubject={Hướng dẫn cài đặt từng bước cho người mới bắt đầu},
  pdfkeywords={Windows, Laravel 12, Vue 3, MySQL, MongoDB, Redis, StayGo}
}

% Cấu hình Header và Footer (chỉ để lại số trang ở footer)
\pagestyle{fancy}
\fancyhf{}
\fancyfoot[C]{\thepage}
\renewcommand{\headrulewidth}{0pt}
\renewcommand{\footrulewidth}{0pt}

% Định dạng tiêu đề các mục
\titleformat{\section}
  {\bfseries\color{PrimaryNavy}\fontsize{13.5}{16.5}\selectfont}
  {\thesection.}{0.5em}{}
\titlespacing*{\section}{0pt}{10pt}{3pt}

\titleformat{\subsection}
  {\bfseries\color{AccentBlue}\fontsize{11}{13.5}\selectfont}
  {\thesubsection}{0.5em}{}
\titlespacing*{\subsection}{0pt}{7pt}{2pt}

% Cấu hình TColorBox
% 1. Hộp lệnh Terminal đẹp mắt (Dành riêng để Copy)
\newtcolorbox{terminalcmd}[2][]{
  enhanced,
  breakable,
  colback=TerminalBg,
  colframe=TerminalBorder,
  arc=2.5mm,
  boxrule=0.8pt,
  left=4mm,right=4mm,top=2.5mm,bottom=2.5mm,
  fonttitle=\sffamily\bfseries\small, title={\textcolor{white}{[LỆNH COPY CHẠY] -- }\textcolor{TerminalPrompt}{#2}},
  coltitle=white,
  attach boxed title to top left={yshift=-2.2mm,xshift=4mm},
  boxed title style={boxrule=0pt,colframe=AccentBlue,colback=AccentBlue,arc=1.2mm},
  fontupper=\ttfamily\color{TerminalText}\small,
  #1
}

% 2. Hộp từng bước (Step Box)
\newtcolorbox{stepbox}[2][]{
  enhanced,
  breakable,
  colback=SoftBlueBg,
  colframe=BorderBlue,
  arc=2mm,
  boxrule=0.7pt,
  left=4mm,right=4mm,top=2.5mm,bottom=2.5mm,
  title={\bfseries\color{PrimaryNavy}#2},
  coltitle=PrimaryNavy,
  attach boxed title to top left={yshift=-2mm,xshift=3mm},
  boxed title style={boxrule=0.5pt,colframe=BorderBlue,colback=white,arc=1.2mm},
  #1
}

% 3. Hộp Lưu ý (Note Box)
\newtcolorbox{notebox}[1][LƯU Ý QUAN TRỌNG]{
  enhanced,
  breakable,
  colback=NoteBg,
  colframe=NoteBorder,
  arc=2mm,
  boxrule=0.7pt,
  left=3.5mm,right=3.5mm,top=2mm,bottom=2mm,
  title={\bfseries\color{NoteText}[GHI CHÚ] #1},
  coltitle=NoteText,
  fontupper=\color{NoteText}\small
}

% 4. Hộp Cảnh báo / Khắc phục lỗi (Warning Box)
\newtcolorbox{warningbox}[1][CẢNH BÁO / KHẮC PHỤC SỰ CỐ]{
  enhanced,
  breakable,
  colback=WarnBg,
  colframe=WarnBorder,
  arc=2mm,
  boxrule=0.7pt,
  left=3.5mm,right=3.5mm,top=2mm,bottom=2mm,
  title={\bfseries\color{WarnText}[CẢNH BÁO] #1},
  coltitle=WarnText,
  fontupper=\color{WarnText}\small
}

% 5. Hộp Mẹo thực hành (Tip Box)
\newtcolorbox{tipbox}[1][MẸO THỰC HÀNH NHANH]{
  enhanced,
  breakable,
  colback=TipBg,
  colframe=TipBorder,
  arc=2mm,
  boxrule=0.7pt,
  left=3.5mm,right=3.5mm,top=2mm,bottom=2mm,
  title={\bfseries\color{TipText}[MẸO HAY] #1},
  coltitle=TipText,
  fontupper=\color{TipText}\small
}

\newcommand{\codeinline}[1]{\colorbox{InlineCodeBg}{\textcolor{AccentBlue}{\texttt{#1}}}}

\begin{document}

% ==================== BÌA TÀI LIỆU TIÊU CHUẨN ====================
\begin{center}
  \vspace*{1mm}
  {\fontsize{17}{21}\selectfont\bfseries\color{PrimaryNavy} HƯỚNG DẪN CÀI ĐẶT VÀ KHỞI CHẠY DỰ ÁN\\TRÊN HỆ ĐIỀU HÀNH WINDOWS\par}
  \vspace{2.5mm}
  \noindent\rule{0.6\textwidth}{0.8pt}
  \vspace{2mm}
\end{center}

\begin{tcolorbox}[colback=white,colframe=BorderBlue,arc=2mm,boxrule=0.7pt,left=4mm,right=4mm,top=2.5mm,bottom=2.5mm]
\small
\textbf{Tên dự án:} Hệ thống Đặt phòng và Quản lý Khách sạn StayGo\\
\textbf{Kiến trúc:} SPA Client-Server + Realtime Event-Driven + Polyglot Persistence\\
\textbf{Công nghệ cốt lõi:} Laravel 12, Vue 3, Node.js 22, MySQL 8.0, MongoDB 8.0, Redis 7\\
\textbf{Mục tiêu tài liệu:} Giúp bất kỳ ai đều có thể tự tay cài đặt, khởi chạy hoàn chỉnh và trải nghiệm toàn bộ tính năng của hệ thống trên máy tính Windows chỉ trong \textbf{5 -- 10 phút}.
\end{tcolorbox}

\vspace{1mm}

% ==================== BƯỚC 1: TẢI MÃ NGUỒN ====================
\section{BƯỚC 1: Tải mã nguồn dự án bằng Git Clone}

\begin{stepbox}{Các thao tác tải mã nguồn và chuyển vào thư mục dự án}
\begin{enumerate}[leftmargin=5mm,itemsep=2pt]
  \item Mở cửa sổ dòng lệnh (\textbf{Command Prompt} hoặc \textbf{PowerShell}) trên máy tính tại thư mục bạn muốn lưu dự án (ví dụ ổ \texttt{D:\textbackslash} hoặc Desktop).
  \item Sao chép và chạy lệnh sau để tải toàn bộ mã nguồn từ GitHub về máy:
\begin{terminalcmd}{Tải mã nguồn từ GitHub về máy tính}
git clone https://github.com/hongphuoc6104/quanlykhachsan.git
\end{terminalcmd}
  \item Sau khi tải xong, chuyển dòng lệnh vào ngay thư mục dự án:
\begin{terminalcmd}{Di chuyển vào thư mục dự án}
cd quanlykhachsan
\end{terminalcmd}
\end{enumerate}
\end{stepbox}

\begin{tipbox}[Mẹo mở lại cửa sổ dòng lệnh tại thư mục dự án nếu lỡ tắt]
Nếu bạn đã lỡ tắt cửa sổ dòng lệnh: Hãy mở thư mục \texttt{quanlykhachsan} bằng \textbf{File Explorer} $\rightarrow$ Nhấp chuột trái vào \textbf{thanh địa chỉ đường dẫn} ở phía trên cùng $\rightarrow$ Gõ chữ \codeinline{cmd} rồi nhấn phím \textbf{Enter}. Cửa sổ dòng lệnh sẽ tự động mở ra ngay tại thư mục dự án.
\end{tipbox}

% ==================== BƯỚC 2: KIỂM TRA FILE CẤU HÌNH .ENV ====================
\section{BƯỚC 2: Chuẩn bị file cấu hình môi trường (.env)}

Ngay tại cửa sổ dòng lệnh trong thư mục \texttt{quanlykhachsan}, hãy tạo file \texttt{.env} từ file mẫu \texttt{.env.example}:

\begin{terminalcmd}{Tạo file .env từ file mẫu trên Windows (chạy trong CMD)}
copy .env.example .env
\end{terminalcmd}

\begin{terminalcmd}{Nếu bạn dùng PowerShell, chạy lệnh này để copy:}
Copy-Item .env.example .env
\end{terminalcmd}

\begin{notebox}[Các cổng kết nối mạng mặc định của hệ thống]
\begin{itemize}[leftmargin=5mm,itemsep=1pt]
  \item \textbf{Cổng 3000:} Dành cho giao diện Web người dùng (Frontend Vue 3).
  \item \textbf{Cổng 8000:} Dành cho máy chủ xử lý dữ liệu (Backend Laravel 12 API).
  \item \textbf{Cổng 3001:} Dành cho máy chủ thời gian thực (Realtime WebSocket Socket.IO).
  \item \textbf{Cổng 3306:} Dành cho Cơ sở dữ liệu MySQL 8.0.
  \item \textbf{Cổng 27017:} Dành cho Cơ sở dữ liệu MongoDB 8.0.
\end{itemize}
\textit{Ghi chú: Nếu máy tính bạn đã cài sẵn MySQL (như XAMPP/WampServer) đang chiếm cổng 3306, hãy mở file \texttt{.env} và đổi thành \codeinline{MYSQL\_HOST\_PORT=3307} hoặc tắt phần mềm XAMPP trước khi khởi chạy.}
\end{notebox}

% ==================== BƯỚC 3: KHỞI CHẠY DỰ ÁN ====================
\section{BƯỚC 3: Khởi chạy toàn bộ hệ thống}

Bạn có thể lựa chọn 1 trong 2 cách sau để khởi chạy:

\subsection*{Cách 1: Click đúp chuột vào file \texttt{run.bat} (Tự động hóa hoàn toàn)}
\begin{enumerate}[leftmargin=5mm,itemsep=1pt]
  \item Mở thư mục dự án \texttt{quanlykhachsan} trên máy tính.
  \item Tìm file có tên \textbf{\texttt{run.bat}} và \textbf{nhấp đúp chuột trái (Double-click)} vào file này.
  \item Cửa sổ sẽ tự động thực hiện toàn bộ việc cấu hình, chuẩn bị cơ sở dữ liệu và khởi chạy các dịch vụ!
\end{enumerate}

\subsection*{Cách 2: Chạy lệnh từ cửa sổ dòng lệnh}
Tại cửa sổ CMD / PowerShell ở thư mục dự án, sao chép và dán lệnh sau:

\begin{terminalcmd}{Khởi chạy toàn bộ hệ thống}
run.bat
\end{terminalcmd}

\begin{terminalcmd}{Hoặc chạy lệnh compose:}
docker compose up -d --build
\end{terminalcmd}

\begin{stepbox}{Hệ thống đang tự động làm những gì trong lần đầu chạy?}
Khi bạn khởi chạy lần đầu tiên, hãy kiên nhẫn đợi từ \textbf{2 đến 5 phút} (tùy vào tốc độ mạng). Hệ thống sẽ tự động thực hiện hoàn toàn các công việc sau:
\begin{enumerate}[leftmargin=5mm,itemsep=1pt]
  \item Tự động xây dựng môi trường cho Backend Laravel 12 và Frontend Vue 3.
  \item Khởi tạo các dịch vụ CSDL MySQL 8.0, MongoDB 8.0 (Replica Set \texttt{rs0}) và Redis 7.
  \item Tự động chạy lệnh tạo bảng CSDL (\textit{Migrations}) và nạp sẵn toàn bộ dữ liệu mẫu (\textit{Seeders}): Danh sách khách sạn, hình ảnh, phòng mẫu, dịch vụ, voucher và tài khoản demo.
  \item Khi bạn thấy thông báo hoàn tất hoặc Frontend báo \texttt{Local: http://localhost:3000} là hệ thống đã hoàn toàn sẵn sàng!
\end{enumerate}
\end{stepbox}

% ==================== BƯỚC 4: TRUY CẬP VÀ TRẢI NGHIỆM ====================
\section{BƯỚC 4: Truy cập và trải nghiệm các chức năng của hệ thống}

Mở trình duyệt web của bạn (Google Chrome, Microsoft Edge, Firefox,...) và truy cập vào các địa chỉ sau:

\begin{table}[htbp]
  \centering
  \fontsize{9}{11.5}\selectfont
  \begin{tabularx}{\textwidth}{>{\bfseries\color{AccentBlue}}p{32mm}>{\ttfamily\color{PrimaryNavy}}p{52mm}>{\raggedright\arraybackslash}X}
    \toprule
    \rowcolor{TableHeader}
    \textcolor{white}{\textbf{Thành phần}} & \textcolor{white}{\textbf{Địa chỉ URL truy cập}} & \textcolor{white}{\textbf{Mô tả chức năng}} \\
    \midrule
    \rowcolor{TableSubHeader}
    Giao diện Web chính & http://localhost:3000 & Trang chủ đặt phòng dành cho khách hàng và đăng nhập quản trị \\
    Cổng Backend API & http://localhost:8000/api/v1 & Cổng giao tiếp REST API v1 cho toàn bộ hệ thống \\
    Kiểm tra Backend & http://localhost:8000/up & Trang theo dõi trạng thái hoạt động của Laravel \\
    Kiểm tra Realtime & http://localhost:3001/health & Trang theo dõi trạng thái kết nối WebSocket Socket.IO \\
    \bottomrule
  \end{tabularx}
\end{table}

\subsection{Danh sách tài khoản dùng thử (Đã tích hợp sẵn phân quyền RBAC)}

Hệ thống đã có sẵn 5 tài khoản mẫu tương ứng với 5 vai trò khác nhau trong quy trình vận hành khách sạn. Bạn hãy đăng nhập thử từng tài khoản tại trang \url{http://localhost:3000/login} để trải nghiệm:

\begin{table}[htbp]
  \centering
  \fontsize{8.5}{11}\selectfont
  \begin{tabularx}{\textwidth}{>{\bfseries}p{22mm}>{\ttfamily}p{44mm}>{\ttfamily}p{26mm}>{\raggedright\arraybackslash}X}
    \toprule
    \rowcolor{TableHeader}
    \textcolor{white}{\textbf{Vai trò (Role)}} & \textcolor{white}{\textbf{Email đăng nhập}} & \textcolor{white}{\textbf{Mật khẩu}} & \textcolor{white}{\textbf{Tính năng kiểm thử nổi bật}} \\
    \midrule
    \rowcolor{white}
    Super Admin & admin@gmail.com & admin123 & Toàn quyền hệ thống, quản lý tất cả khách sạn, nhân sự, cấu hình hệ thống. \\
    \rowcolor{TableSubHeader}
    Quản lý & manager@gmail.com & manager123 & Quản lý khách sạn "An Nhiên Đà Lạt", sửa giá phòng, thêm voucher, dịch vụ, xem báo cáo doanh thu. \\
    \rowcolor{white}
    Lễ tân & receptionist@gmail.com & receptionist123 & \textbf{Xem sơ đồ phòng trực quan Realtime}, xử lý Check-in/Check-out, tạo đơn tại quầy, xuất hóa đơn. \\
    \rowcolor{TableSubHeader}
    Kế toán & accountant@gmail.com & accountant123 & Quản lý danh sách hóa đơn, đối soát dòng tiền thanh toán, kiểm tra hoàn tiền khi hủy phòng. \\
    \rowcolor{white}
    Khách hàng & customer@gmail.com & customer123 & Tìm kiếm phòng theo ngày, đặt phòng, chọn voucher, thanh toán giả lập (VietQR/Card), chat CSKH realtime. \\
    \bottomrule
  \end{tabularx}
\end{table}

% ==================== BƯỚC 5: CÁC CÂU LỆNH THƯỜNG DÙNG ====================
\section{BƯỚC 5: Các câu lệnh quản trị và vận hành thường dùng}

Mỗi khi cần thao tác bảo trì, hãy mở cửa sổ CMD / PowerShell tại thư mục dự án và dùng các lệnh sau:

\begin{terminalcmd}{Xem danh sách và trạng thái hoạt động của các dịch vụ}
docker compose ps
\end{terminalcmd}

\begin{terminalcmd}{Xem nhật ký (Logs) thời gian thực của hệ thống khi cần kiểm tra lỗi}
docker compose logs -f
\end{terminalcmd}

\begin{terminalcmd}{Tạm dừng và tắt toàn bộ hệ thống một cách an toàn}
docker compose down
\end{terminalcmd}

\begin{terminalcmd}{Bật lại hệ thống khi đã tải xong lần trước (khởi động cực nhanh)}
docker compose up -d
\end{terminalcmd}

\begin{terminalcmd}{Khôi phục / Nạp lại toàn bộ dữ liệu mẫu (Seeder) về trạng thái ban đầu}
docker compose exec backend php artisan db:seed --force
\end{terminalcmd}

\begin{terminalcmd}{Chạy toàn bộ bộ kiểm thử tự động (Unit Test / Feature Test) của Backend}
docker compose exec backend php artisan test
\end{terminalcmd}

% ==================== BƯỚC 6: XỬ LÝ LỖI THƯỜNG GẶP ====================
\section{BƯỚC 6: Hướng dẫn xử lý các sự cố thường gặp trên Windows}

\begin{warningbox}[1. Lỗi: "Bind for 0.0.0.0:3306 failed: port is already allocated"]
\textbf{Nguyên nhân:} Cổng 3306 trên máy tính của bạn đã bị phần mềm khác chiếm dụng (thường do cài XAMPP, WampServer, hoặc dịch vụ MySQL cài trực tiếp trên Windows).\\
\textbf{Cách xử lý 1:} Tắt phần mềm XAMPP / MySQL Service trên Windows.\\
\textbf{Cách xử lý 2:} Mở file \texttt{.env} trong thư mục dự án, tìm dòng \texttt{MYSQL\_HOST\_PORT=3306} và đổi thành \codeinline{MYSQL\_HOST\_PORT=3307}, sau đó chạy lại lệnh \codeinline{docker compose up -d}.
\end{warningbox}

\begin{warningbox}[2. Lỗi: "Bind for 0.0.0.0:27017 failed: port is already allocated"]
\textbf{Nguyên nhân:} Cổng 27017 bị trùng do máy tính đã cài sẵn MongoDB Service.\\
\textbf{Cách xử lý:} Mở file \texttt{.env}, tìm dòng \texttt{MONGODB\_HOST\_PORT=27017} và đổi thành \codeinline{MONGODB\_HOST\_PORT=27018}, sau đó chạy lại lệnh.
\end{warningbox}

\begin{warningbox}[3. Lỗi: Trang web Frontend hiện màn hình trắng hoặc báo không kết nối được API]
\textbf{Nguyên nhân:} Trình duyệt chưa tải kịp hoặc Backend container đang trong quá trình nạp dữ liệu migration.\\
\textbf{Cách xử lý:} Đợi khoảng 30 giây rồi nhấn \codeinline{Ctrl + F5} trên trình duyệt để tải lại trang. Đảm bảo file \texttt{.env} có cấu hình \codeinline{VITE\_API\_BASE\_URL=http://localhost:8000/api/v1}.
\end{warningbox}

\vspace{3mm}
\begin{center}
  \noindent\rule{0.6\textwidth}{0.8pt}\\
  \vspace{1.5mm}
  {\small\color{gray}Chúc bạn cài đặt thành công và có trải nghiệm tuyệt vời với hệ thống StayGo!}
\end{center}

\end{document}
